\cleardoublepage

{
    \sectionnonum{附录}

    \appendixsubsecmajornumbering
    \subsection{12 工具的完整接口规范}
    
    \autoref{tab:appendix-tool-schemas}
    对第四章 \autoref{tab:tool-list} 的 12 个工具补充给出完整的输入参数与关键输出字段说明，
    便于读者按字段名复现 DAGPlanner 中的
    \verb|${sX.field}| 占位符。
    
    \small
    
    \begin{longtable}{
    p{0.18\linewidth}
    p{0.30\linewidth}
    p{0.30\linewidth}
    p{0.14\linewidth}
    }
    
    \caption{\label{tab:appendix-tool-schemas}
    12 工具的完整接口规范（输入参数 + 输出字段）}
    \\
    
    \toprule
    
    \textbf{工具名}
    &
    \textbf{关键输入参数}
    &
    \textbf{关键输出字段}
    &
    \textbf{用途}
    
    \\
    
    \midrule
    \endfirsthead
    
    \toprule
    
    \textbf{工具名}
    &
    \textbf{关键输入参数}
    &
    \textbf{关键输出字段}
    &
    \textbf{用途}
    
    \\
    
    \midrule
    \endhead
    
    \midrule
    \multicolumn{4}{r}{续下页}
    \\
    \endfoot
    
    \bottomrule
    \endlastfoot
    
    \texttt{matrix\_operation}
    &
    \texttt{operation}
    (add/sub/mul/inverse/det),
    \texttt{matrix\_a},
    \texttt{matrix\_b}
    &
    \texttt{result},
    \texttt{shape},
    \texttt{determinant},
    \texttt{real},
    \texttt{imag}
    &
    矩阵基本运算与行列式
    \\
    
    \texttt{numerical}
    &
    \texttt{expression},
    \texttt{lower},
    \texttt{upper},
    \texttt{variable}
    &
    \texttt{integral},
    \texttt{absolute\_error},
    \texttt{expression},
    \texttt{interval}
    &
    一元函数自适应数值积分
    \\
    
    \texttt{curve\_fitting}
    &
    \texttt{model}
    (linear/poly2/exponential/power),
    \texttt{x},
    \texttt{y}
    &
    \texttt{model},
    \texttt{params},
    \texttt{r\_squared},
    \texttt{formula}
    &
    多模型曲线拟合
    \\
    
    \texttt{equation\_solver}
    &
    \texttt{expression},
    \texttt{lower},
    \texttt{upper},
    \texttt{variable}
    &
    \texttt{root},
    \texttt{residual},
    \texttt{expression},
    \texttt{interval}
    &
    区间内方程求根
    \\
    
    \texttt{statistics}
    &
    \texttt{values} (数组)
    &
    \texttt{count},
    \texttt{mean},
    \texttt{median},
    \texttt{std},
    \texttt{variance},
    \texttt{min},
    \texttt{max},
    \texttt{range},
    \texttt{q1},
    \texttt{q3},
    \texttt{iqr}
    &
    描述性统计量
    \\
    
    \texttt{hypothesis\_test}
    &
    \texttt{test\_type}
    (one\_sample\_t /
    two\_sample\_t),
    \texttt{sample\_a},
    \texttt{sample\_b},
    \texttt{alpha}
    &
    \texttt{test\_type},
    \texttt{t\_statistic},
    \texttt{p\_value},
    \texttt{alpha},
    \texttt{reject\_null}
    &
    t 检验等假设检验
    \\
    
    \texttt{linear\_regression}
    &
    \texttt{x},
    \texttt{y}
    &
    \texttt{slope},
    \texttt{intercept},
    \texttt{r\_squared},
    \texttt{p\_value},
    \texttt{std\_err},
    \texttt{formula}
    &
    一元线性回归
    \\
    
    \texttt{correlation}
    &
    \texttt{x},
    \texttt{y},
    \texttt{method}
    (pearson/spearman)
    &
    \texttt{method},
    \texttt{correlation},
    \texttt{p\_value},
    \texttt{strength},
    \texttt{direction}
    &
    相关性分析
    \\
    
    \texttt{line\_chart}
    &
    \texttt{x},
    \texttt{y},
    \texttt{title},
    \texttt{xlabel},
    \texttt{ylabel},
    \texttt{output\_path}
    &
    \texttt{file\_path},
    \texttt{width\_pixels},
    \texttt{height\_pixels}
    &
    折线图
    \\
    
    \texttt{scatter\_chart}
    &
    \texttt{x},
    \texttt{y},
    \texttt{title},
    \texttt{xlabel},
    \texttt{ylabel},
    \texttt{output\_path}
    &
    \texttt{file\_path},
    \texttt{width\_pixels},
    \texttt{height\_pixels}
    &
    散点图
    \\
    
    \texttt{bar\_chart}
    &
    \texttt{labels},
    \texttt{values},
    \texttt{title},
    \texttt{xlabel},
    \texttt{ylabel},
    \texttt{output\_path}
    &
    \texttt{file\_path},
    \texttt{width\_pixels},
    \texttt{height\_pixels}
    &
    条形图
    \\
    
    \texttt{heatmap}
    &
    \texttt{matrix},
    \texttt{row\_labels},
    \texttt{col\_labels},
    \texttt{title},
    \texttt{output\_path}
    &
    \texttt{file\_path},
    \texttt{width\_pixels},
    \texttt{height\_pixels}
    &
    热力图
    \\
    
    \end{longtable}
    
    \normalsize
    ```
    

    \subsection{O1--O4 提示工程规则全文}

    第四章 \autoref{tab:o1-o4} 给出了 4 条规则的概要，本节给出其在 SYSTEM\_PROMPT 中的实际中文措辞，便于读者复用或修订。

    \subsubsection{O1: 结果合理性检查}

    \begin{quote}\small
    （从 observation 取值前必须执行。）
    拟合类工具返回的 $R^2$：若 $<0$ 或 $>1$，视为工具输出异常。
    拟合参数的量级：若拟合参数与输入数据主导量级相差超过 2 个数量级（例如数据在 $10^3$ 量级却拟合出 $10^0$ 量级的常数），视为异常。
    返回值中出现 NaN / Inf / 空值，视为异常。
    一旦检测到上述异常，下一轮 \emph{必须} 尝试不同的解决路径；可选方案包括：换一个拟合模型（例如从 exponential 改为 linear / polynomial）；对数据做变换再拟合；换一个工具解决同一目标。
    在异常被纠正之前，不得将异常工具返回的数值写入最终答复。
    \end{quote}

    \subsubsection{O2: 数值计算规范}

    \begin{quote}\small
    涉及 3 个或以上样本点的统计量（求和、均值、中位数、方差、标准差、四分位数、分位数）\emph{必须} 通过 \texttt{descriptive\_statistics} 工具计算。严禁在 thought / content 里以任何形式写出 “(a+b+c+...)/n”、“mean = ...”、“sum = ...” 等手算表达式。
    线性代数运算（矩阵求逆、乘积、行列式、特征值、转置后再乘等）\emph{必须} 通过 \texttt{matrix\_operation}。
    积分、方程求根、曲线拟合 \emph{必须} 通过对应专业工具。
    最终答复报告数字时，至少保留工具返回原值的 4--6 位有效数字。
    \end{quote}

    \subsubsection{O3: 开场规划模板}

    \begin{quote}\small
    对于需要 3 步及以上的任务，在第一轮回复的 content 中 \emph{必须} 先以下面的格式列出计划：

    \texttt{Plan:}\\
    \texttt{1. 用 <工具名> 做 <子目标>，得到 <中间变量>}\\
    \texttt{2. 用 <工具名> 做 <子目标>，用到 <中间变量>}\\
    \texttt{...}

    写完 Plan 后，同一轮直接发起第一个 tool\_call。
    后续每一轮继续推进 Plan：在 \texttt{思考：} 前缀里简要引用当前走到 Plan 的第几步。
    若在执行中发现 Plan 有偏差，明确修正 Plan 后再执行。
    \end{quote}

    \subsubsection{O4: Few-shot 示例注入}

    O4 在 SYSTEM\_PROMPT 末尾追加一条与本题集无关、用于演示正确工作风格的完整范例（异常切换 + 最终答复格式）。
    范例本身较长，完整内容见 \texttt{src/prompts/system\_prompts.py} 中的 \texttt{FEWSHOT\_EXAMPLE} 常量。

    \subsection{ext 压力测试用例清单}

    第五章测试集中的 8 道 \texttt{ext\_*} 压力题专门用于检验 RATS 与 DAGPlanner 的算法贡献，其设计目标见\autoref{tab:appendix-ext-cases}。

    \begin{table}[ht]
        \centering
        \caption{\label{tab:appendix-ext-cases}8 道 \texttt{ext\_*} 压力测试用例的设计目标}
        \begin{tabularx}{\linewidth}{clXc}
            \toprule
            编号 & 用例 ID & 设计目标 & 主要类别 \\ \midrule
            1 & \texttt{ext\_rs\_01} & 在 \texttt{descriptive\_statistics} 与 \texttt{hypothesis\_test} 之间区分（统计描述 vs 假设检验） & statistics \\
            2 & \texttt{ext\_rs\_02} & 在 \texttt{linear\_regression} 与 \texttt{correlation\_analysis} 之间区分（拟合 vs 相关） & statistics \\
            3 & \texttt{ext\_rs\_03} & 在 \texttt{equation\_solver} 与 \texttt{numerical\_integration} 之间区分 & numerical \\
            4 & \texttt{ext\_rs\_04} & 在 \texttt{heatmap} 与 \texttt{bar\_chart} 之间区分（多维 vs 一维可视化） & visualization \\
            5 & \texttt{ext\_dag\_01} & 5 步链式：行列式 $\to$ 求根 $\to$ 积分 $\to$ 统计 $\to$ 折线图 & composite \\
            6 & \texttt{ext\_dag\_02} & 4 步链式：统计 $\to$ 相关 $\to$ 回归 $\to$ 散点图 & composite \\
            7 & \texttt{ext\_dag\_03} & 3 步链式 + 中间结果分支：拟合 + 残差统计 + 残差直方图 & composite \\
            8 & \texttt{ext\_dag\_04} & 4 步链式：矩阵特征值 $\to$ 排序 $\to$ 选取 $\to$ 条形图 & composite \\ \bottomrule
        \end{tabularx}
    \end{table}

    \removeappendixsubsecmajornumbering
}
